%%%%%%%%%%%%%%%%%%%%%%%%%%%%%%%%%%%%%%%%%%%%%%%%%%%%%%%%%%%%%%%%%%%%%
%  sub/kr-1-introduction.tex  —  1. 서론 (한글본)
%  영문본 sub/en-1-introduction.tex 의 번역본. 구조를 동일하게 유지한다.
%%%%%%%%%%%%%%%%%%%%%%%%%%%%%%%%%%%%%%%%%%%%%%%%%%%%%%%%%%%%%%%%%%%%%
\section{서론}\label{sec:introduction-kr}

% 서론에서는 연구가 다루는 과학교육 문제와 그 중요성을 제시하고, 선행연구를
% 비판적으로 검토하여 연구 공백(gap)을 드러낸 뒤, 본 연구의 목적과 연구 문제를
% 명확히 밝힙니다.

여기에 서론을 작성합니다. 연구가 다루는 과학교육 문제와 그 중요성을 밝히고,
선행연구를 검토하여 연구 공백을 드러낸 뒤, 명시적인 연구 문제로 마무리합니다.

% 인용은 \citep{key}(괄호 인용) 또는 \citet{key}(본문 인용)로 씁니다.
% 예: 선행연구는 이 문제를 다루었으며 \citep{example2024}, \citet{example2024}는
% \dots 라고 논하였다.
